\section{Sparse Visual Features and Their Summaries}
\label{sec:surrogates}
The construction needs nonnegative feature descriptions and a declared
object set. It does not depend on a particular neural architecture.
We separate the feature extractor, the sparse surrogate, and the spatial
summary so that each operation has an explicit interpretation.

\subsection{Activation states and sparse surrogates}
\label{sub:activation-states}
Let $X$ be a finite nonempty image collection. For image $x\in X$, let
$P_{x}$ be a finite nonempty set of spatial sites. A frozen vision model
provides an activation $h_{x,p}\in\R^{m}$ at site $p\in P_{x}$ of a fixed
layer. In a convolutional network these sites index a feature map; in a
vision transformer they can index patch tokens. Classification and
register tokens must be handled separately if the intended interpretation
is spatial.

A sparse autoencoder (SAE) encodes and reconstructs the same activation:
\begin{equation}\label{eq:sae}
  z_{x,p}=E(h_{x,p})\in\R_{\geq 0}^{d},\qquad
  \widehat h_{x,p}=D(z_{x,p}),\qquad
  D(z)=W_{\mathrm{dec}}z+b_{\mathrm{dec}}.
\end{equation}
Typically $d>m$. A nonnegative Top-$k$ encoder takes a rectified linear
projection and retains its $k$ largest entries, so its code has at most
$k$ nonzero coordinates~\cite{gao2024}. With this hard sparsity constraint,
one can minimize squared reconstruction error. An alternative is a
rectified encoder with an $\ell_{1}$ penalty:
\begin{equation}\label{eq:sae-loss}
  \mathcal{L}_{\mathrm{SAE}}
  =\mathbb{E}_{h}\bigl[
    \|D(E(h))-h\|_{2}^{2}+\lambda\|E(h)\|_{1}\bigr].
\end{equation}
The pilot uses the hard Top-$k$ constraint, not the penalized objective.
Decoder normalization removes a simple source of arbitrary scaling;
training and random initialization can still change the dictionary.

A transcoder instead takes the input $a_{x,p}$ to a selected transformation
and predicts its output $h_{x,p}$:
\begin{equation}\label{eq:transcoder}
  z_{x,p}=E_{\mathrm{TC}}(a_{x,p}),\qquad
  \widehat h_{x,p}=D_{\mathrm{TC}}(z_{x,p}).
\end{equation}
It is a sparse surrogate for a computation rather than merely for an
activation state~\cite{dunefsky2024transcoders}. Both models fit the theory
when their codes are nonnegative. Their coordinates are not directly
comparable across dictionaries, even when dimensions coincide. The pilot
below evaluates SAEs only; no SAE--transcoder advantage is inferred.

\subsection{Images, regions, and spatial sites}
\label{sub:spatial-objects}
For an image-level analysis, define its summary by
\begin{equation}\label{eq:max-summary}
  u_{x,j}=\max_{p\in P_{x}}z_{x,p,j},\qquad 1\leq j\leq d.
\end{equation}
This is an existential summary for each coordinate separately. It retains
activation magnitudes but discards which site supplies each maximum.
A region-level analysis can instead use objects $(x,S)$ with nonempty
$S\subseteq P_{x}$ and take the maximum over $S$. A site-level analysis
uses objects $(x,p)$ and assigns the unpooled code $z_{x,p}$. All three
choices yield the same kind of graded context; they answer different
questions.

For a fixed image context, let
\begin{equation}\label{eq:vector-lattice}
  M_{j}=\max_{x\in X}u_{x,j},\qquad
  \U=\prod_{j=1}^{d}[0,M_{j}],\qquad \X=\mathcal{P}(X).
\end{equation}
Order $\U$ coordinatewise and $\X$ by inclusion. A complete lattice is an
ordered set in which every family, including the empty family, has a
least upper bound (join) and a greatest lower bound (meet). In $\U$ these
operations are coordinatewise suprema and infima, with bottom $0$ and
top $M$. In $\X$ they are unions and intersections, with bottom
$\varnothing$ and top $X$. A coordinate with $M_{j}=0$ simply contributes a
one-element factor. Sparsity of site codes need not survive max pooling:
different sites can activate different coordinates.

The bound $M$ is specific to the context on which closure is computed.
For held-out retrieval, a frozen training query is compared directly with
new codes without clipping them at $M$. Thus the experiment does not
silently use test activations to choose training thresholds. A larger
context can be defined afterwards, but its closure operator can differ.

\subsection{Available vision libraries}
\label{sub:libraries}
Public implementations make the framework applicable beyond the pilot.
ViT-Prisma supplies activation access, vision SAEs, and CLIP transcoders;
its published resources include CLIP and DINO dictionaries
\cite{joseph2025prisma,vitprisma}.
SAEV provides a vision SAE training and inference workflow with linked
Hugging Face checkpoints~\cite{saev,stevens2025vision}.
Overcomplete offers several SAE and dictionary-learning variants together
with visualization and evaluation tools~\cite{overcomplete}.
PatchSAE provides closely related evidence that sparse visual codes can
support patch-level analysis during model adaptation
\cite{lim2025patchsae}.

For a foundation-model extension, ViT-Prisma is the first candidate when
both SAEs and transcoders are required; SAEV is a natural alternative for
vision SAE experiments, and Overcomplete for controlled training studies.
This is a choice of software scope, not a measured performance ranking.
A checkpoint is usable only with its matching backbone, layer, activation
normalization, preprocessing, and token selection. The implementation
therefore accepts cached nonnegative site codes as its core interface.
The executed pilot uses the project's existing PyTorch environment and a
small explicitly defined surrogate, rather than claiming to have tested
these external libraries.
